Scientific applications execute on heterogeneous systems with multiple accelerators and deep memory and storage hierarchies, where data may traverse GPU memory, CPU memory, node-local storage, burst buffers, and parallel file systems. The volumes are substantial. A single output step of a trillion-particle Vector Particle-In-Cell (VPIC) plasma simulation writes tens of terabytes~\cite{byna2012trillion}. As data moves across these tiers, applications increasingly need computation applied along the data path: compression, filtering, format conversion, encryption, and feature extraction for downstream analysis~\cite{patwary2015bdcats, tan2022analysis}. Performing this work while data is in motion reduces storage pressure and movement cost, improves security, and prepares data for consumption without a separate post-processing stage. Efficient and transparent processing along the data path is therefore becoming a core requirement of HPC data management systems, yet the parallel file systems that anchor these systems expose a read/write interface and orchestrate nothing.

The need to compute along the data path is not new, and the HPC community has pursued it from two directions that differ in \emph{where} computation runs. \emph{In-situ} processing executes on the same nodes as the simulation, sharing its memory or address space, and so avoids moving data at the cost of competing with the simulation for resources. \emph{In-transit} processing ships data to separate staging resources and computes there as the data travels toward storage, trading movement cost for isolation. Systems in both categories have made substantial progress toward a common goal---moving computation closer to data. Existing in-situ systems such as DataSpaces~\cite{dataspaces}, FlexPath~\cite{flexpath}, and ADIOS~\cite{godoy2020adios2} support tight coupling, staging, and derived data computation alongside running simulations. In-transit approaches such as Runway~\cite{ravi2023runway} and the derived data work of Gainaru et al.~\cite{gainaru2024toderive} make runtime-level decisions about when and where to execute computation as data moves.

These systems share three limitations. First, the operations are invoked explicitly by the application rather than declared once and applied thereafter: Runway requires PDC API calls before each region transfer~\cite{ravi2023runway}, so the transformation schedule is written into the application and must be updated application code whenever the pipeline or the hardware changes. Second, the operations are expressed as isolated steps rather than composable pipelines. HDF5 chains filters within a single dataset's write path~\cite{hdf5}, but the chain is fixed at dataset creation and cannot branch on data state or runtime conditions; Runway schedules individual transformations with no abstraction for chaining them at all. Third, support is confined to one kind of operation, or splits the two kinds across mechanisms that do not interoperate. HDF5 filters transform byte streams and cannot derive new datasets, while ADIOS~2 supports derived variables~\cite{gainaru2024toderive} and per-variable compression operators~\cite{godoy2020adios2} through separate interfaces, so attaching a transformation to a derived result is not expressible.

We use \emph{in-flight processing} for an approach that cuts across this distinction. It is defined not by where computation runs but by when: during the transfer itself, with the location (e.g., simulation node, staging resource, or destination) selected by the runtime rather than fixed by the programming model. A system built this way can place an operation where it is cheapest at the moment it executes, which neither an in-situ nor an in-transit programming model expresses.

Building a general in-flight processing system raises five challenges. \textbf{(C1) Execution must be transparent to the application.} If an application names an operation before each transfer, it retains the orchestration the system is meant to absorb, and the four challenges below become its problem rather than the runtime's. Pipelines must instead be declared once and carried out by ordinary reads and writes. \textbf{(C2) Operations depend on data state.} Data that has already been compressed must be handled differently from data that has not, and different portions of a dataset may exist in different states at the same point in a workflow. Analysis compounds this: the system must know which derived objects have been computed from which inputs, and whether a derived result should be persisted or discarded once consumed. The system needs a representation that tracks state and identifies valid paths through it. \textbf{(C3) Operations have dependencies that span objects.} An analysis may require several independently produced inputs before it can run and may emit several outputs, each bound for a different destination. Resolving readiness and routing results cannot be left to the application without reintroducing exactly the orchestration burden the system is meant to remove. \textbf{(C4) Placement depends on runtime conditions.} The best device for a given operation varies with runtime conditions such as device utilization, operation cost, and data movement overhead. Consequently, a static assignment to either a CPU or GPU cannot consistently achieve efficient execution. \textbf{(C5) Operations must not stall data movement.} Simulations of this kind run for thousands of timesteps, each involving substantial computation, which leaves room to overlap output with ongoing work. Applying computation synchronously during I/O forfeits that overlap, so operations must be integrated with in-flight transfers rather than layered on top of them.

We present Data Flyway, a runtime that meets these challenges by representing pipelines as directed graphs registered with I/O resources. We implement it in Proactive Data Containers (PDC)~\cite{tang2018toward, pdcdocs}, an object-centric HPC data management system whose servers sit between the application and persistent storage. At up to 128 nodes, Data Flyway hides transformation cost behind application compute on a write-heavy checkpointing workload, and sustains roughly 55\% higher I/O throughput on a deep learning training workload that re-reads its dataset each epoch. Selecting devices at runtime rather than statically reduces total transformation time by 15\% at 2048 parallel MPI ranks.

\paragraph{Contributions.} This paper makes the following contributions.
\begin{itemize}
    \item A declarative model in which an application states the data state it wants rather than invoking operations. A pipeline is described once, registered with the runtime, and attached to an object or region, after which ordinary reads and writes drive it and orchestration leaves application code entirely (\S\ref{sec:pipeline}).

    \item A directed-graph representation that unifies transformation and analysis in one structure. Nodes are data states and named objects, edges are operations, and the runtime tracks which state each region holds and which derived objects already exist.

    \item A runtime scheduler that selects a device for each operation from predicted cost rather than a fixed assignment. Its GPU cost model augments payload size and device utilization with lag features taken from the preceding invocation, capturing transfer autocorrelation that a size-and-utilization model does not (\S\ref{sec:model}, \S\ref{sec:sched}).

    \item An implementation in PDC and an evaluation at up to 128 nodes across checkpointing, deep learning training, and multi-stage analysis workloads.
\end{itemize}
